% ============================================================================
% Fichier  : partie0/P00_paiement_28_decembre.tex
% Titre    : Prologue -- Le Paiement du 28 Décembre
% Auteur   : MINKA MI NGUIDJOÏ Thierry Emmanuel
% Version  : 0.1 -- novembre 2026
% ============================================================================
% Faits du dossier MINÉQUIR à garder identiques dans tout le cours :
%  - marché n° 0412/M/MINÉQUIR/2025, BATIRURAL SARL (Mbalmayo), pistes rurales ;
%    décompte n° 3, facture F-2025-117, 187 450 000 FCFA ;
%  - mandat 2025-08817 (22/12/2025, circuit normal), payé le 24/12 sur le
%    compte n° 1 (banque A) ;
%  - 27/12/2025 21 h 52 : RIB modifié par le compte mbarga.j (compte n° 2,
%    banque B, au nom de « BATIRURAL SERVICES ») ;
%  - 28/12/2025 23 h 14 : mandat 2025-09342, facture « F-2025-117 bis »,
%    liquidé et ordonnancé par mbarga.j ; payé le 30/12 ;
%  - MBARGA, ex-chef du service de la liquidation, muté à Ebolowa le
%    14/08/2025 ; compte jamais désactivé ;
%  - janvier 2026 : inspection sur pièces, 40 dossiers papier de décembre,
%    « régularité satisfaisante » ;
%  - 16/03/2026 : ordre de mission du CONSUPE (vérificatrice NGO BIYONG) ;
%    extraction de 9 611 mandats 2025 ; 3 paires de doublons, dont 1 seule
%    avec RIB différent ;
%  - lecture des 8 contrôles attendus : 2 OUI, 5 NON, 1 NEK.
% Personnages : EBOUMBOU (DAF), ATANGANA (chef de la cellule informatique),
% FOTSO (chef du poste comptable), NGO BIYONG (vérificatrice du CONSUPE).
% ============================================================================

\chapter[Prologue : Le Paiement du 28 Décembre]{Prologue : Le Paiement du 28 Décembre}
\label{chap:prologue}
\chaptermark{Prologue}

\epigraph{%
    « Le dossier était en ordre~; le système, non. »%
}{Note de synthèse d'une vérificatrice}

\niveauChapitre{Introductif}{%
    Aucun prérequis informatique. Il suffit de connaître le circuit de la
    dépense publique~: engagement, liquidation, ordonnancement, paiement.
    Ce prologue se lit avant toute définition~: il part d'un contrôle qui
    n'a rien vu.%
}

\begin{boxobjectifs}
\begin{itemize}
    \item Reconstituer, à partir des pièces et des traces informatiques, la
          chronologie d'un double paiement (dossier MINÉQUIR).
    \item Expliquer pourquoi un contrôle sur pièces, bien conduit, peut ne
          pas voir une irrégularité qui se trouve dans le système.
    \item Distinguer une constatation, appuyée sur une preuve conservée, de
          l'analyse que l'on en tire.
    \item Identifier les contrôles internes qui auraient dû empêcher ou
          détecter le double paiement, et les coter \OUI{}, \NON{} ou
          \NEK{}.
    \item Citer le texte qui donne au Contrôle Supérieur de l'État la
          compétence du contrôle des systèmes d'information.
    \item Énoncer les trois réflexes de l'auditeur public face à un système
          d'information~: le mandat, la preuve, le contradictoire.
\end{itemize}
\end{boxobjectifs}

\bigskip

% ============================================================================
\section*{Avant de commencer : un contrôle qui n'a rien vu}
% ============================================================================

Yaoundé, lundi 16~mars~2026. Mme~NGO~BIYONG, vérificatrice au Contrôle
Supérieur de l'État, reçoit son ordre de mission~: vérifier l'exécution du
budget 2025 de la Direction des Affaires Financières du ministère de
l'Équipement rural, le MINÉQUIR\footnote{Le MINÉQUIR, ses agents, ses
fournisseurs, les dates et les montants de ce dossier sont fictifs. Les
mécanismes du circuit de la dépense, les contrôles et les textes cités sont
exacts.}. Dans le dossier permanent, une pièce attire son attention~: le
rapport d'une inspection conduite en janvier sur les dépenses de décembre.
Il examine quarante dossiers de paiement et conclut à une «~régularité
satisfaisante~».

Trois semaines plus tard, la même vérificatrice établit qu'une facture a
été payée deux fois en décembre, la seconde fois sur un compte bancaire qui
n'appartient pas au fournisseur. Le montant en jeu est de
\textbf{187\,450\,000~francs CFA}.
L'inspection de janvier ne s'était pas trompée sur les dossiers qu'elle a
examinés~: ils étaient complets et réguliers. Elle n'a simplement jamais vu
le second paiement, parce qu'il n'avait pas de dossier.

\medskip

Ce prologue reconstitue ce dossier. Il montre ce qu'un contrôle sur pièces
peut voir et ce qu'il ne peut pas voir, ce que la trace informatique
révèle, et quels contrôles internes ont manqué. Il en tire les trois
réflexes qui structurent tout le cours. Les notions y sont nommées~; elles
sont définies avec précision dans les chapitres suivants.

\begin{boxalert}
\textbf{Violation typique.} Tirer l'échantillon d'un contrôle dans les
classeurs de dossiers papier, alors que les paiements sont produits par une
application~: ce qui n'a pas de dossier ne peut pas être tiré.

\smallskip
\textbf{Règle~:} la population d'un contrôle est l'ensemble des opérations
enregistrées dans le système, pas l'ensemble des dossiers rangés.
\end{boxalert}

\separateur

% ============================================================================
\section{Ce qui s'est passé}
\label{sec:pro-chronologie}
% ============================================================================

\subsection{Le marché et le circuit}

Le MINÉQUIR a passé en 2025 le marché n°~0412/M/MINÉQUIR/2025 avec la
société BATIRURAL~SARL, de Mbalmayo, pour l'aménagement de pistes rurales.
Le décompte n°~3 de ce marché, facturé sous le numéro F-2025-117, s'élève
à 187\,450\,000~FCFA.

Depuis 2024, la dépense du ministère passe par une application, SIDÉPENSE.
Chaque étape du circuit y est un écran et un profil d'utilisateur~:
l'engagement réserve le crédit, la liquidation constate le service fait et
arrête le montant, l'ordonnancement émet le mandat, puis le poste comptable
reçoit le mandat par voie électronique et paie. Les pièces justificatives
sont scannées et jointes au mandat~; le dossier papier est classé à la
Direction des Affaires Financières.

\begin{tableaucours}{L{4.0cm}L{8.2cm}}{%
    \tableauHeader{Élément} &
    \tableauHeader{Description}
}
Entité & Direction des Affaires Financières (DAF) du MINÉQUIR, Yaoundé \\
Application & SIDÉPENSE~: engagement, liquidation, ordonnancement,
transmission des mandats au poste comptable, fichier des fournisseurs \\
Volume 2025 & 9\,611 mandats \\
Acteurs & M.~EBOUMBOU, directeur des affaires financières~; M.~ATANGANA,
chef de la cellule informatique~; M.~FOTSO, chef du poste comptable \\
Contrôle & Inspection sur pièces (janvier 2026)~; mission du CONSUPE
(mars 2026), conduite par Mme~NGO~BIYONG \\
\end{tableaucours}
\vspace{-0.8em}
\begin{center}{\small\itshape Tableau~P.1. Le dossier MINÉQUIR}\end{center}

\subsection{Reconstituer avant de juger}

Une vérificatrice ne commence pas par chercher un responsable~: elle
reconstitue la suite des événements, et rattache chaque événement à une
pièce. Ce qu'on lui raconte sans pièce reste une déclaration, à vérifier.
Le Tableau~P.2 donne la chronologie établie par Mme~NGO~BIYONG.

\begin{tableaucours}{L{2.6cm}L{6.4cm}L{3.2cm}}{%
    \tableauHeader{Date} &
    \tableauHeader{Événement} &
    \tableauHeader{Pièce}
}
14 août 2025 & M.~MBARGA, chef du service de la liquidation, est muté à
Ebolowa & Décision de mutation \\
22 déc. 2025 & Mandat 2025-08817 (décompte n°~3) émis par le circuit
normal & Mandat, dossier papier complet \\
24 déc. 2025 & Paiement sur le compte n°~1 de BATIRURAL (banque~A) &
Relevé du poste comptable \\
27 déc. 2025, 21~h~52 & Le RIB de BATIRURAL est remplacé, dans le fichier
des fournisseurs, par un compte n°~2 (banque~B) au nom de «~BATIRURAL
SERVICES~» & Journal de SIDÉPENSE \\
28 déc. 2025, 23~h~14 & Mandat 2025-09342, facture «~F-2025-117 bis~»,
même montant, liquidé et ordonnancé par le compte \texttt{mbarga.j} &
Journal de SIDÉPENSE, mandat électronique \\
30 déc. 2025 & Paiement du second mandat sur le compte n°~2 & Relevé du
poste comptable \\
Janv. 2026 & Inspection de 40 dossiers papier de décembre~: «~régularité
satisfaisante~» & Rapport d'inspection \\
16 mars 2026 & Ordre de mission du CONSUPE & Ordre de mission \\
23 mars 2026 & Extraction des 9\,611 mandats de 2025 et du journal des
opérations & Fichiers extraits, empreintes \\
\end{tableaucours}
\vspace{-0.8em}
\begin{center}{\small\itshape Tableau~P.2. Chronologie du double paiement
(reconstitution de la vérificatrice)}\end{center}

Ce tableau ne contient que des faits. Il ne dit pas encore qui a agi, ni
pourquoi~: le compte \texttt{mbarga.j} a servi, mais rien n'établit à ce
stade que M.~MBARGA l'ait utilisé lui-même. Séparer ce qui est établi de ce
que l'on en déduit est la première discipline de l'auditeur.

\begin{boxdef}
\termecle{Constatation}\index{constatation} (\textit{finding})~: fait
vérifiable, établi par l'auditeur et appuyé sur une preuve conservée. Un
tiers compétent qui refait la même démarche sur les mêmes pièces aboutit au
même résultat.

\medskip
\termecle{Analyse}\index{analyse}~: interprétation que l'auditeur tire
d'une ou plusieurs constatations, au regard d'un critère, avec un niveau de
confiance déclaré (faible, moyen ou élevé).
\end{boxdef}

\separateur

% ============================================================================
\section{Ce que le contrôle sur pièces a vu}
\label{sec:pro-pieces}
% ============================================================================

L'inspection de janvier a fait correctement ce qu'elle a fait. Elle a tiré
quarante dossiers dans les classeurs de décembre, vérifié pour chacun
l'engagement, le service fait, la facture, le visa et l'imputation. Les
quarante dossiers étaient complets. Le mandat 2025-08817 en faisait partie.

Son erreur n'est pas dans l'examen, elle est dans la \textbf{population}.
L'échantillon a été tiré dans les classeurs, c'est-à-dire parmi les
mandats qui avaient un dossier papier. Le mandat 2025-09342 n'en avait pas~:
il a été créé directement dans l'application, avec un scan de la facture du
premier mandat, renommée «~bis~». Aucun tirage dans les classeurs ne
pouvait le faire sortir.

\begin{boxscenario}{Dossier MINÉQUIR. Deux populations}
\textbf{Constatation.} Les classeurs de décembre contiennent 1\,184
dossiers~; l'application enregistre 1\,191 mandats émis en décembre. Sept
mandats n'ont pas de dossier papier, dont le mandat 2025-09342.

\medskip
\textbf{Analyse} (confiance élevée). L'inspection a conclu sur la
population des dossiers rangés, pas sur celle des mandats émis. Sa
conclusion vaut pour les 1\,184 dossiers~; elle ne dit rien des sept
autres. Un rapprochement entre les deux populations, fait avant le tirage,
aurait suffi à les faire apparaître.
\end{boxscenario}

Cette différence entre deux populations est la première leçon du
cours~: dès que l'administration s'informatise, le dossier papier n'est
plus le système, il n'en est qu'une sortie.

\separateur

% ============================================================================
\section{Ce que la trace informatique révèle}
\label{sec:pro-trace}
% ============================================================================

\subsection{L'extraction}

Mme~NGO~BIYONG ne part pas des classeurs, mais de l'application. Elle
demande par écrit à M.~ATANGANA une extraction de tous les mandats de 2025
et du journal des opérations, assiste à l'extraction, calcule l'empreinte
des fichiers et la fait contresigner. Elle vérifie ensuite que le total des
mandats extraits est égal au total des paiements du compte administratif.
Le Chapitre~\ref{chap:preuve} détaille ces précautions~: sans elles, une
extraction n'est qu'un fichier remis par l'audité.

\subsection{Le test des doublons}

Sur les 9\,611 mandats, elle recherche ceux qui portent le même
fournisseur et le même montant. Le test, fait dans un tableur, renvoie
trois paires. Deux s'expliquent par des pièces (deux livraisons identiques
d'un même fournisseur). La troisième est celle de BATIRURAL~: même
fournisseur, même montant, même numéro de facture au suffixe près, et deux
comptes bancaires différents.

\begin{boxoutils}{Recherche des doublons dans un tableur}
Dans la feuille des mandats (colonnes~: B fournisseur, D montant), une
colonne de clé et une colonne de comptage suffisent~:
\begin{lstlisting}
F2 : =B2&"|"&D2
G2 : =NB.SI($F$2:$F$9612;F2)
\end{lstlisting}
Un filtre sur \texttt{G~>~1} affiche les mandats qui partagent le même
fournisseur et le même montant. Le test ne prouve rien à lui seul~: il
désigne les mandats à examiner, pièce par pièce.
\end{boxoutils}

\subsection{Le journal}

Le journal de SIDÉPENSE enregistre, pour chaque opération sensible,
l'utilisateur, la date, l'heure et l'action. Il permet de reconstituer
les deux opérations du Tableau~P.3.

\begin{lstlisting}[basicstyle=\ttfamily\scriptsize]
2025-12-27 21:52:07 mbarga.j FOURNISSEUR MODIF_RIB   BATIRURAL  cpte1->cpte2
2025-12-28 23:09:41 mbarga.j MANDAT      LIQUIDATION 2025-09342 187450000
2025-12-28 23:14:18 mbarga.j MANDAT      ORDONNANCE  2025-09342 187450000
2025-12-30 10:02:55 pc.fotso MANDAT      PAIEMENT    2025-09342 cpte2
\end{lstlisting}
\vspace{-0.6em}
\begin{center}{\small\itshape Tableau~P.3. Extrait du journal de SIDÉPENSE
(fichier extrait le 23~mars~2026, empreinte au dossier)}\end{center}

\begin{boxscenario}{Dossier MINÉQUIR. Ce que dit le journal}
\textbf{Constatations.}
\begin{enumerate}[(a)]
    \item Le compte \texttt{mbarga.j} était actif le 28~décembre~2025,
          quatre mois après la mutation de son titulaire.
    \item Le même compte a modifié un RIB, liquidé et ordonnancé le même
          mandat~; aucune deuxième personne n'est intervenue.
    \item Les deux opérations ont eu lieu un samedi soir et un dimanche
          soir, hors des heures de service.
    \item SIDÉPENSE a accepté un second mandat de même fournisseur, même
          montant et même numéro de facture au suffixe «~bis~» près.
\end{enumerate}

\medskip
\textbf{Analyse} (confiance moyenne). Le double paiement n'est pas une
erreur de saisie~: la modification du RIB la veille, sur le seul compte
du fournisseur concerné, et l'heure des opérations désignent une
manœuvre. L'identité de la personne qui a utilisé le compte n'est pas
établie~: un compte resté actif peut être utilisé par quelqu'un qui en
connaît le mot de passe.
\end{boxscenario}

\separateur

% ============================================================================
\section{Les contrôles internes qui ont manqué}
\label{sec:pro-controles}
% ============================================================================

La question de l'auditeur n'est pas seulement «~qui a fait cela~?~», qui
relève d'autres procédures, mais «~qu'est-ce qui aurait dû l'empêcher, ou
le détecter, et pourquoi cela n'a-t-il pas fonctionné~?~». C'est la
question du \termecle{contrôle interne}\index{contrôle interne}~: l'ensemble
des dispositifs qui donnent une assurance raisonnable que les opérations
sont régulières, les informations fiables et les biens protégés. Le
Chapitre~\ref{chap:sci} le définit selon le COSO, et le
Chapitre~\ref{chap:eval-sci} apprend à l'évaluer.

Mme~NGO~BIYONG dresse la liste des contrôles qu'elle s'attendait à trouver
sur ce circuit, et cote chacun. Elle répond \OUI{} si le contrôle existe et
a fonctionné, \NON{} s'il manque ou n'a pas fonctionné, \NEK{}
(\textit{Not Enough Knowledge}) si elle n'a pas encore de quoi conclure.

\begin{tableaucours}{L{0.7cm}L{6.6cm}L{1.2cm}L{3.6cm}}{%
    \tableauHeader{N°} &
    \tableauHeader{Contrôle attendu} &
    \tableauHeader{Cote} &
    \tableauHeader{Preuve}
}
1 & Les comptes des agents mutés sont désactivés à la date de la mutation
& \NON{} & Liste des comptes actifs, décision de mutation \\
2 & Un même profil ne peut pas liquider et ordonnancer & \NON{} & Matrice
des profils de SIDÉPENSE \\
3 & Toute modification de RIB exige une seconde validation & \NON{} &
Journal~: une seule intervention \\
4 & L'application bloque un mandat de même fournisseur, même montant et
même facture & \NON{} & Mandat 2025-09342 accepté \\
5 & Le poste comptable vérifie les pièces jointes avant paiement & \NEK{} &
Pièces du poste comptable demandées, non reçues \\
6 & Les opérations sensibles sont journalisées & \OUI{} & Journal
extrait \\
7 & Les paiements de fin de gestion font l'objet d'une revue & \NON{} &
Aucune revue produite \\
8 & Les journaux sont conservés et protégés & \OUI{} & Journal intact
depuis 2024 \\
\end{tableaucours}
\vspace{-0.8em}
\begin{center}{\small\itshape Tableau~P.4. Les contrôles attendus sur le
circuit de la dépense}\end{center}

\begin{boxlecture}{Circuit de la dépense, MINÉQUIR}
\textbf{Constatation.} Sur huit contrôles attendus, deux existent et
fonctionnent, cinq manquent ou n'ont pas fonctionné, un ne peut pas encore
être coté.

\medskip
\textbf{Analyse} (confiance élevée). Les deux contrôles qui fonctionnent
sont des contrôles de \emph{détection} a posteriori~: ils ont permis de
reconstituer les faits, pas de les empêcher. Tous les contrôles de
\emph{prévention} ont manqué. La cote \NEK{} du contrôle~5 n'est pas une
faiblesse de la mission~: elle dit honnêtement ce qui reste à établir, et
désigne la prochaine pièce à obtenir.

\medskip
\hfill\conformite{2}{8}{1}
\end{boxlecture}

\begin{boxinfo}
\textbf{Pourquoi \NEK{} et pas \NON{}~?} Si le poste comptable devait
produire une trace de sa vérification et ne l'a pas produite, la réponse
est \NON{}~: l'absence d'une trace exigée est une constatation. Ici, la
vérificatrice n'a pas encore reçu les pièces qu'elle a demandées~: elle ne
sait pas. Elle l'écrit, et relance.
\end{boxinfo}

\separateur

% ============================================================================
\section{Qui peut contrôler ce système}
\label{sec:pro-mandat}
% ============================================================================

Mme~NGO~BIYONG n'a pas eu à demander la permission d'auditer SIDÉPENSE.
Le texte organique du Contrôle Supérieur de l'État le lui permet
expressément.

\begin{boxjuris}
\textbf{\loicam{Décret n°~2013/287 du 4~septembre~2013} portant
organisation des services du Contrôle supérieur de l'État, article~2,
alinéa~3}~\cite{decret2013287consupe}~: \textit{« Dans l'exercice de leurs
attributions, les services du Contrôle supérieur de l'Etat effectuent~:
un contrôle de conformité et de régularité~; un contrôle financier~; un
contrôle de performance~; un contrôle de l'environnement~; un contrôle des
systèmes d'information. »}

\medskip
\textbf{Article~2, alinéa~4}~: \textit{« Les contrôles sont effectués par
les vérificateurs du Contrôle supérieur de l'Etat sur pièces et sur place.
A cet effet, et nonobstant les contrôles effectués par les autres organes
de contrôle, les responsables des structures auditées sont tenus de
présenter les pièces relatives à la période auditée. »}

\medskip
\textbf{Portée pour l'audit.} Le contrôle des systèmes d'information est
une compétence écrite du CONSUPE, au même titre que le contrôle de
régularité. L'obligation de présenter les pièces couvre les pièces
électroniques~: une extraction, un journal, une matrice des profils sont
des pièces relatives à la période auditée.
\end{boxjuris}

Deux autres institutions regardent les systèmes d'information publics,
chacune sous un angle différent. L'\sigle{ANTIC} conduit les audits de
sécurité des réseaux et des systèmes d'information prévus par la
\loicam{loi n°~2010/012}. La Chambre des Comptes de la Cour suprême juge
les comptes des comptables publics~; elle rencontre les systèmes
d'information par extension, parce que ces comptes en sont issus. Le
Chapitre~\ref{chap:institutions} précise qui fait quoi, et comment ces
interventions s'articulent.

\separateur

% ============================================================================
\section{Trois réflexes}
\label{sec:pro-reflexes}
% ============================================================================

Le dossier MINÉQUIR fait apparaître les trois réflexes que le cours
installe, chapitre après chapitre.

\begin{enumerate}
    \item \textbf{Le mandat.} Savoir au nom de quoi l'on contrôle, sur quel
          périmètre et selon quels critères. C'est l'ordre de mission qui
          a permis d'exiger l'extraction~; c'est le critère («~un mandat
          ne peut pas être payé deux fois~») qui a permis de coter
          (Chapitres~\ref{chap:institutions} et~\ref{chap:preparer}).
    \item \textbf{La preuve.} Partir du système, pas du classeur~;
          obtenir la donnée soi-même ou y assister, en garantir
          l'intégrité, et conserver tout ce qui fonde une cote
          (Chapitre~\ref{chap:preuve}).
    \item \textbf{Le contradictoire.} Soumettre les constatations à
          l'entité avant de conclure. Le directeur des affaires
          financières et le chef de la cellule informatique recevront le
          Tableau~P.4 et pourront le contester pièce par pièce
          (Chapitre~\ref{chap:constater}).
\end{enumerate}

\begin{boxprocedure}
\textbf{Le contenu minimal d'un papier de travail}, pour chaque contrôle
coté~:
\begin{itemize}
    \item le contrôle attendu et le critère qui le fonde~;
    \item le test effectué, la date et la personne qui l'a fait~;
    \item la pièce obtenue (nom du fichier, empreinte, source)~;
    \item la cote (\OUI{}, \NON{} ou \NEK{}) et la constatation qui la
          justifie~;
    \item pour un \NEK{}, la pièce manquante et la relance.
\end{itemize}
\end{boxprocedure}

\begin{boxmetier}{Vérificateur}
Avant tout tirage d'échantillon, rapprochez le nombre et le montant des
opérations du système avec ceux de la source que vous allez échantillonner
(classeurs, bordereaux). Un écart, même de quelques unités, désigne les
opérations qu'il faut examiner en premier.
\end{boxmetier}

\begin{boxmetier}{Inspecteur}
Un rapport qui conclut «~régularité satisfaisante~» doit dire sur quelle
population il porte. Écrivez-le~: «~sur les 1\,184 dossiers classés de
décembre~», et non «~sur les dépenses de décembre~».
\end{boxmetier}

\begin{boxmetier}{Contrôleur financier}
Les paiements de fin de gestion concentrent le risque~: volume élevé,
délais courts, personnel réduit. Une revue des mandats émis hors des
heures de service et des RIB modifiés dans les derniers jours de
l'exercice est un contrôle simple, et il aurait arrêté le second paiement.
\end{boxmetier}

\separateur

% ============================================================================
\section{Le chemin du cours}
\label{sec:pro-chemin}
% ============================================================================

Le cours se déroule en huit journées de quatre heures. Les deux dernières
heures de la huitième journée sont offertes.

\begin{tableaucours}{L{1.3cm}L{8.6cm}L{2.3cm}}{%
    \tableauHeader{Jour} &
    \tableauHeader{Question traitée} &
    \tableauHeader{Chapitres}
}
J01 & Qu'est-ce qu'auditer un système d'information public~? & P00,
\ref{chap:audit-public} \\
J02 & Qu'est-ce qu'un système de contrôle interne, et comment
l'évaluer~? & \ref{chap:sci}, \ref{chap:eval-sci} \\
J03 & Qui contrôle quoi, et selon quelles normes~? &
\ref{chap:institutions}, \ref{chap:normes} \\
J04 & Comment préparer la mission et obtenir la preuve~? &
\ref{chap:preparer}, \ref{chap:preuve} \\
J05 & Comment conclure~? Que vaut le contrôle des accès~? &
\ref{chap:constater}, \ref{chap:acces} \\
J06 & Que valent l'exploitation et les contrôles des applications~? &
\ref{chap:exploitation}, \ref{chap:applications} \\
J07 & Que disent les données~? Que valent les projets et les
prestataires~? & \ref{chap:donnees-fraude}, \ref{chap:gouvernance} \\
J08 & Une mission complète, en équipe & \ref{chap:cas-integrateur} \\
\end{tableaucours}
\vspace{-0.8em}
\begin{center}{\small\itshape Tableau~P.5. Les huit journées}\end{center}

Trois dossiers accompagnent le cours. Le MINÉQUIR revient dans les
chapitres sur le contrôle interne, la mission, les accès et les
applications. La \emph{CAPNO}, caisse publique de prévoyance de Garoua,
sert aux chapitres sur la preuve et l'analyse des données~: son fichier de
paie contient des agents qui n'existent pas. Le projet \emph{e-RECETTES},
à Bafoussam, sert aux chapitres sur l'exploitation, les projets et au cas
intégrateur.

\separateur

\begin{boxresume}
\begin{itemize}
    \item \textbf{Le dossier n'est pas le système.} Dans une administration
          informatisée, le classeur n'est qu'une sortie de l'application~;
          la population d'un contrôle se prend dans le système.
    \item \textbf{Reconstituer avant de juger.} Chaque événement est
          rattaché à une pièce~; ce qui n'a pas de pièce reste une
          déclaration.
    \item \textbf{Constatation et analyse se séparent.} La constatation est
          vérifiable par un tiers~; l'analyse porte un niveau de
          confiance.
    \item \textbf{La question du contrôle interne.} Ce qui aurait dû
          empêcher ou détecter l'irrégularité, et pourquoi cela n'a pas
          fonctionné.
    \item \textbf{Trois réponses.} \OUI{}, \NON{}, \NEK{}~; l'absence d'une
          trace exigée vaut \NON{}, une pièce non encore obtenue vaut
          \NEK{}.
    \item \textbf{Une compétence écrite.} Le décret n°~2013/287 range le
          contrôle des systèmes d'information parmi les cinq contrôles du
          CONSUPE.
    \item \textbf{Trois réflexes.} Le mandat, la preuve, le contradictoire.
\end{itemize}
\end{boxresume}

\begin{boxexo}[title={\ding{45}\quad Exercice~P.1. Constatation ou analyse~? \niveaubase}]
Classez chacun des énoncés suivants et justifiez en une ligne.
\begin{enumerate}[(a)]
    \item «~Le compte \texttt{mbarga.j} a ordonnancé le mandat 2025-09342
          le 28~décembre~2025 à 23~h~14 (journal extrait le 23~mars,
          empreinte au dossier).~»
    \item «~M.~MBARGA a détourné 187\,450\,000~FCFA.~»
    \item «~Sept mandats de décembre n'ont pas de dossier papier.~»
    \item «~Le contrôle des modifications de RIB n'est pas efficace.~»
    \item «~La cellule informatique semble débordée.~»
\end{enumerate}
\end{boxexo}

\begin{boxexo}[title={\ding{45}\quad Exercice~P.2. Coter un contrôle \niveaubase}]
Pour chacune des situations suivantes, donnez la cote (\OUI{}, \NON{} ou
\NEK{}) et justifiez.
\begin{enumerate}[(a)]
    \item La procédure exige un procès-verbal mensuel de revue des comptes
          actifs~; la DAF n'en produit aucun pour 2025.
    \item Le chef de la cellule informatique affirme que les mots de passe
          expirent tous les 90~jours~; le paramétrage n'a pas encore été
          montré.
    \item La matrice des profils, extraite par la vérificatrice, ne donne
          à aucun profil à la fois la liquidation et l'ordonnancement.
    \item Le test de doublon a été paramétré en juin 2026, après les
          faits.
\end{enumerate}
\end{boxexo}

\begin{boxexo}[title={\ding{45}\quad Exercice~P.3. Deux populations \niveaumoyen}]
L'inspection de janvier a examiné 40 dossiers parmi 1\,184.
\begin{enumerate}[(a)]
    \item Quelle était la probabilité que le mandat 2025-09342 figure dans
          son échantillon~? Justifiez.
    \item Proposez, en trois étapes, la démarche qui aurait permis de le
          détecter sans augmenter la taille de l'échantillon.
    \item Rédigez la phrase de conclusion que l'inspection aurait dû
          écrire, compte tenu de sa population réelle.
\end{enumerate}
\end{boxexo}

\begin{boxexo}[title={\ding{45}\quad Exercice~P.4. Le premier constat \niveauexpert}]
À partir du Tableau~P.4 et du journal (Tableau~P.3)~:
\begin{enumerate}[(a)]
    \item Rédigez, en séparant constatation et analyse, le constat portant
          sur le contrôle~2 (séparation entre liquidation et
          ordonnancement).
    \item Indiquez le critère que vous citez, et la pièce qui fonde la
          cote.
    \item Formulez une recommandation adressée au directeur des affaires
          financières, et une adressée au chef de la cellule
          informatique.
    \item Dites quelles réponses vous attendez de chacun pendant la
          procédure contradictoire, et ce qui pourrait vous faire changer
          la cote.
\end{enumerate}
\end{boxexo}

\begin{boxinfo}
\textbf{Transition.} Le prologue a montré un contrôle qui n'a rien vu et
une vérificatrice qui a vu, parce qu'elle est partie du système. Le
Chapitre~\ref{chap:audit-public} pose le vocabulaire~: ce que signifient
contrôler, auditer, vérifier, et ce qui change quand l'objet du contrôle
est un système d'information public.
\end{boxinfo}
